% norme di progetto
% da compilare con il comando pdflatex Norme_di_progetto_v_x.x.x.tex

% Dichiarazioni di ambiente e inclusione di pacchetti
% da usare tramite il comando % Dichiarazioni di ambiente e inclusione di pacchetti
% da usare tramite il comando % Dichiarazioni di ambiente e inclusione di pacchetti
% da usare tramite il comando \input{../../util/hx-ambiente}

\documentclass[a4paper,titlepage]{article}
\usepackage[T1]{fontenc}
\usepackage[utf8]{inputenc}
\usepackage[english,italian]{babel}
\usepackage{microtype}
\usepackage{lmodern}
\usepackage{underscore}
\usepackage{graphicx}
\usepackage{eurosym}
\usepackage{float}
\usepackage{fancyhdr}
\usepackage[table]{xcolor}
\usepackage{longtable}
\usepackage{chngpage}
\usepackage{grffile}
\usepackage[titles]{tocloft}
\usepackage{hyperref}
\hypersetup{hidelinks}

\usepackage{../../util/hx-vers}
\usepackage{../../util/hx-macro}
\usepackage{../../util/hx-front}

% solo se si vuole una nuova pagina ad ogni \section:
\usepackage{titlesec}
\newcommand{\sectionbreak}{\clearpage}

% stile di pagina:
\pagestyle{fancy}

% solo se si vuole eliminare l'indentazione ad ogni paragrafo:
\setlength{\parindent}{0pt}

% intestazione:
\lhead{\Large{\proj}}
\rhead{\includegraphics[keepaspectratio=true,width=50px]{../../util/hivex_logo2.png}}
\renewcommand{\headrulewidth}{0.4pt}

% pie' di pagina:
\lfoot{\email}
\rfoot{\thepage}
\cfoot{}
\renewcommand{\footrulewidth}{0.4pt}

% spazio verticale tra le celle di una tabella:
\renewcommand{\arraystretch}{1.5}

% profondità di indicizzazione:
\setcounter{tocdepth}{4}
\setcounter{secnumdepth}{4}


\documentclass[a4paper,titlepage]{article}
\usepackage[T1]{fontenc}
\usepackage[utf8]{inputenc}
\usepackage[english,italian]{babel}
\usepackage{microtype}
\usepackage{lmodern}
\usepackage{underscore}
\usepackage{graphicx}
\usepackage{eurosym}
\usepackage{float}
\usepackage{fancyhdr}
\usepackage[table]{xcolor}
\usepackage{longtable}
\usepackage{chngpage}
\usepackage{grffile}
\usepackage[titles]{tocloft}
\usepackage{hyperref}
\hypersetup{hidelinks}

\usepackage{../../util/hx-vers}
\usepackage{../../util/hx-macro}
\usepackage{../../util/hx-front}

% solo se si vuole una nuova pagina ad ogni \section:
\usepackage{titlesec}
\newcommand{\sectionbreak}{\clearpage}

% stile di pagina:
\pagestyle{fancy}

% solo se si vuole eliminare l'indentazione ad ogni paragrafo:
\setlength{\parindent}{0pt}

% intestazione:
\lhead{\Large{\proj}}
\rhead{\includegraphics[keepaspectratio=true,width=50px]{../../util/hivex_logo2.png}}
\renewcommand{\headrulewidth}{0.4pt}

% pie' di pagina:
\lfoot{\email}
\rfoot{\thepage}
\cfoot{}
\renewcommand{\footrulewidth}{0.4pt}

% spazio verticale tra le celle di una tabella:
\renewcommand{\arraystretch}{1.5}

% profondità di indicizzazione:
\setcounter{tocdepth}{4}
\setcounter{secnumdepth}{4}


\documentclass[a4paper,titlepage]{article}
\usepackage[T1]{fontenc}
\usepackage[utf8]{inputenc}
\usepackage[english,italian]{babel}
\usepackage{microtype}
\usepackage{lmodern}
\usepackage{underscore}
\usepackage{graphicx}
\usepackage{eurosym}
\usepackage{float}
\usepackage{fancyhdr}
\usepackage[table]{xcolor}
\usepackage{longtable}
\usepackage{chngpage}
\usepackage{grffile}
\usepackage[titles]{tocloft}
\usepackage{hyperref}
\hypersetup{hidelinks}

\usepackage{../../util/hx-vers}
\usepackage{../../util/hx-macro}
\usepackage{../../util/hx-front}

% solo se si vuole una nuova pagina ad ogni \section:
\usepackage{titlesec}
\newcommand{\sectionbreak}{\clearpage}

% stile di pagina:
\pagestyle{fancy}

% solo se si vuole eliminare l'indentazione ad ogni paragrafo:
\setlength{\parindent}{0pt}

% intestazione:
\lhead{\Large{\proj}}
\rhead{\includegraphics[keepaspectratio=true,width=50px]{../../util/hivex_logo2.png}}
\renewcommand{\headrulewidth}{0.4pt}

% pie' di pagina:
\lfoot{\email}
\rfoot{\thepage}
\cfoot{}
\renewcommand{\footrulewidth}{0.4pt}

% spazio verticale tra le celle di una tabella:
\renewcommand{\arraystretch}{1.5}

% profondità di indicizzazione:
\setcounter{tocdepth}{4}
\setcounter{secnumdepth}{4}


\version{0.0.3}
\creaz{24 dicembre 2016}
\author{\GG, \MM}
\supervisor{\LB, \AZ}
\uso{interno}
\dest{Tutti i membri del gruppo}
\title{Norme di progetto}

\begin{document}
\maketitle
% diario delle modifiche per il piano di progetto
% da includere con % diario delle modifiche per il piano di progetto
% da includere con % diario delle modifiche per il piano di progetto
% da includere con \include{diario}

\begin{diario}
	1.0.1 & {\GG} (Responsabile) & 03/02/2017 & Correzione degli errori individuati nelle prime tre sezioni del documento \\ \hline
	1.0.0 & {\PB} (Responsabile) & 09/01/2017 & Approvazione documento \\ \hline
	0.1.0 & {\MM} (Verificatore) & 07/01/2017 & Verifica documento \\ \hline
	0.0.9 & {\PB} (Responsabile) & 05/01/2017 & Stesura Organigramma \\ \hline
	0.0.8 & {\LB} (Responsabile) & 05/01/2017 & Stesura Consuntivo di Periodo \\ \hline
	0.0.7 & {\LB} (Responsabile) & 04/01/2017 & Stesura Preventivo di Periodo \\ \hline
	0.0.6 & {\LB} (Responsabile) & 03/01/2017 & Inserimento Gantt Diagrammi delle Attività \\ \hline
	0.0.5 & {\PB} (Responsabile) & 29/12/2016 & Stesura Pianificazione delle Attività \\ \hline
	0.0.4 & {\PB} (Responsabile) & 28/12/2016 & Stesura Introduzione \\ \hline
	0.0.3 & {\LB} (Responsabile) & 27/12/2016 & Stesura Modello di sviluppo \\ \hline
	0.0.2 & {\PB} (Responsabile) & 27/12/2016 & Stesura Analisi dei Rischi \\ \hline
	0.0.1 & {\LB} (Responsabile) & 26/12/2016 & Stesura scheletro \\ \hline
\end{diario}

\begin{diario}
	1.0.1 & {\GG} (Responsabile) & 03/02/2017 & Correzione degli errori individuati nelle prime tre sezioni del documento \\ \hline
	1.0.0 & {\PB} (Responsabile) & 09/01/2017 & Approvazione documento \\ \hline
	0.1.0 & {\MM} (Verificatore) & 07/01/2017 & Verifica documento \\ \hline
	0.0.9 & {\PB} (Responsabile) & 05/01/2017 & Stesura Organigramma \\ \hline
	0.0.8 & {\LB} (Responsabile) & 05/01/2017 & Stesura Consuntivo di Periodo \\ \hline
	0.0.7 & {\LB} (Responsabile) & 04/01/2017 & Stesura Preventivo di Periodo \\ \hline
	0.0.6 & {\LB} (Responsabile) & 03/01/2017 & Inserimento Gantt Diagrammi delle Attività \\ \hline
	0.0.5 & {\PB} (Responsabile) & 29/12/2016 & Stesura Pianificazione delle Attività \\ \hline
	0.0.4 & {\PB} (Responsabile) & 28/12/2016 & Stesura Introduzione \\ \hline
	0.0.3 & {\LB} (Responsabile) & 27/12/2016 & Stesura Modello di sviluppo \\ \hline
	0.0.2 & {\PB} (Responsabile) & 27/12/2016 & Stesura Analisi dei Rischi \\ \hline
	0.0.1 & {\LB} (Responsabile) & 26/12/2016 & Stesura scheletro \\ \hline
\end{diario}

\begin{diario}
	1.0.1 & {\GG} (Responsabile) & 03/02/2017 & Correzione degli errori individuati nelle prime tre sezioni del documento \\ \hline
	1.0.0 & {\PB} (Responsabile) & 09/01/2017 & Approvazione documento \\ \hline
	0.1.0 & {\MM} (Verificatore) & 07/01/2017 & Verifica documento \\ \hline
	0.0.9 & {\PB} (Responsabile) & 05/01/2017 & Stesura Organigramma \\ \hline
	0.0.8 & {\LB} (Responsabile) & 05/01/2017 & Stesura Consuntivo di Periodo \\ \hline
	0.0.7 & {\LB} (Responsabile) & 04/01/2017 & Stesura Preventivo di Periodo \\ \hline
	0.0.6 & {\LB} (Responsabile) & 03/01/2017 & Inserimento Gantt Diagrammi delle Attività \\ \hline
	0.0.5 & {\PB} (Responsabile) & 29/12/2016 & Stesura Pianificazione delle Attività \\ \hline
	0.0.4 & {\PB} (Responsabile) & 28/12/2016 & Stesura Introduzione \\ \hline
	0.0.3 & {\LB} (Responsabile) & 27/12/2016 & Stesura Modello di sviluppo \\ \hline
	0.0.2 & {\PB} (Responsabile) & 27/12/2016 & Stesura Analisi dei Rischi \\ \hline
	0.0.1 & {\LB} (Responsabile) & 26/12/2016 & Stesura scheletro \\ \hline
\end{diario}
\tableofcontents

\section{Introduzione}

\subsection{Scopo del documento}
Questo documento (interno al gruppo) regolamenta i processi del progetto didattico e va quindi letto da ciascun membro del gruppo. Le convenzioni qui prescritte servono a:
\begin{itemize}
	\item garantire ordine all'interno dei documenti e delle varie parti della configurazione del prodotto;
	\item mantenere coerenza nelle notazioni e nelle procedure;
	\item minimizzare i conflitti tra i vari ruoli;
	\item garantire che l'infrastruttura di lavoro sia il più possibile semplice e gestibile --- quindi fruibile;
\end{itemize}

\subsection{Scopo del prodotto}
\scopo % da ../../util/hx-macro.sty

	\subsection{Glossario}
	\subsection{Riferimenti}
		\subsubsection{Normativi}
		\begin{itemize}
		\item ISO/IEC 12207-1995: \url{http://www.math.unipd.it/\~tullio/IS-1/2009/Approfondimenti/ISO_12207-1995.pdf}
		\end{itemize}
\section{Processo di sviluppo}

	\subsection{Studio di fattibilità}
	La fase antecedente alla realizzazione del progetto\ped{g} consiste nel fissare delle riunioni in modo che il team\ped{g} possa discutere di esso. Sarà poi compito degli analisti redigere lo \texttt{studio di fattibilità},la cui struttura è specificata a seguire: %inserire struttura studio di fattibilità
	\subsection{Analisi dei requisiti}
	Sarà sempre compito degli analisti al termine dello studio di fattibilità di produrre l'\texttt{analisi dei requisiti}.Questo documento conterrà i requisiti raccolti ed i casi uso individuati nelle riunioni svolte.
	Viene elencata nelle sezioni a seguire la struttura obbligatoria di requisiti e casi d'uso.
		\subsubsection{Requisiti}
		\subsubsection{Casi d'uso}
		Ogni caso d'uso viene identificato nel seguente modo:
		\centerline{UCx}
		Ogni macro caso d'uso sarai poi analizzato in modo più specifico ed identificato da UC[p].[f] dove p indica il il caso d'uso padre e f il figlio(ogni figlio potrà a sua volta avere altri figli).
		Inoltre ogni caso'uso dovrà avere le seguenti proprietà:
		\begin{itemize}
		\item Nome che lo identifica;
		\item Attori coinvolti in esso;
		\item Descrizione di esso;
		\item Precondizioni se presenti;
		\item Postcondizioni se presenti;
		\item Scenario principale che descrive esso tramite i casi d'uso figli;
		\item Inclusioni se specificate;
		\item Estensioni se specificate;
		\item Scenari alternativi che descrive esse tramite casi d'uso non presenti tra i figli.
		\end{itemize}
	\subsection{Progettazione}
	Il compito dei progettisti consiste di progettare l'architettura del software. % da completare.
	
	\subsection{Codifica}
		\subsubsection{Formattazione del codice}
	\subsection{Strumenti}
		\subsubsection{PragmaDB}
		Per tenere traccia dei casi d'uso e dei requisiti il team\ped{g} ha scelto di usare PragmaDB(https://github.com/StefanoMunari/PragmaDB) creato da studenti degli anni passati. Essa è accessibile via broswer nel server dedicato tramite autenticazione.Inoltre permette l'esportazione in \LaTeX{} delle tabelle.
		\subsubsection{Astah}
		\subsubsection{Ide ancora da scegliere}
	\subsection{Tecnologie utilizzate}
	\subsubsection{HTML5}
	\subsubsection{CSS3}
	\subsubsection{Javascript}
	\subsubsection{Java tomcat o node.js}

\section{Processi di supporto}
I processi normati in questa sezione hanno la funzione di aiutare gli altri processi.

\subsection{Processo di documentazione}
Per gestire la documentazione di progetto si è pensato, innanzitutto, di creare un ramo a parte (intitolato \texttt{doc}) a partire dal ramo \texttt{master} del nostro repository; questa scelta è stata dettata dal basso grado di dipendenza tra la documentazione il resto delle parti della configurazione del prodotto. Riportiamo qui di seguito le norme che riguardano lo sviluppo della documentazione; quelle enunciate con il condizionale vanno interpretate come consigli.

\subsubsection{Ambiente di lavoro} I documenti andranno mantenuti in una struttura ad albero ben organizzata ma non troppo profonda; il miglior compromesso tra organizzazione e profondità dell'albero è risultata essere quella in figura [figura...].

\subsubsection{Nomi ed estensioni} Nome ed estensione dei file sono regolati nel seguente modo:
\paragraph{Nome} Il nome di un file o di una directory non dovrebbe contenere spazi, al fine di facilitare operazioni da riga di comando.
\paragraph{Formato} Ogni documento dev'essere in formato \LaTeX{} (esportabile in PDF su foglio A4); la prima versione di un documento può essere in altri formati ma va presto sostituita da una versione in \LaTeX.
\paragraph{Identificazione} Ogni versione di un documento è identificata univocamente dal nome del file, composto di una prima stringa esprimente la sua funzione (ad esempio “Piano\_di\_progetto”) e una seconda stringa con il numero di versionamento (ad esempio “\_0.0.2”).
\paragraph{Numero di versionamento} Il numero di versionamento di un documento si compone di tre numeri: il primo viene incremementato ad ogni \textbf{approvazione} del documento; il secondo ad ogni \textbf{revisione}; il terzo ad ogni \textbf{aggiunta o modifica} sostanziale.

\subsubsection{Frontespizio} Ogni documento, anche se breve, deve avere un frontespizio recante: titolo del documento, versione, nome del nostro gruppo, nome del progetto, data, autori, revisori, uso (interno o esterno), lista di distribuzione. Per fare questo, il redattore deve inserire \texttt{\textbackslash input\{../../util/hx-ambiente\}} all'inizio del documento.

\subsubsection{Diario delle modifiche} Ogni documento deve includere un diario delle modifiche che elenca i cambiamenti apportati ad ogni versione rispetto alla precedente, assieme all'autore e alla data di tali cambiamenti. Per questo, il redattore deve includere un file intitolato \texttt{diario.tex} nella stessa cartella del documento; al suo interno, deve usare le macro \texttt{\textbackslash begin\{diario\}} e \texttt{\textbackslash end\{diario\}}.

\subsubsection{Ciclo di vita di un documento} Per ogni documento creato, i redattori stendono una bozza (in \LaTeX) che dev'essere poi controllata dai verificatori; se questi rilevano errori o possibili miglioramenti da apportare, segnalano il fatto al \gloss{responsabile di progetto} --- che provvede a rendicontare le ore di lavoro aggiuntive --- e modificano il documento dopo aver avvisato i redattori di tale documento.

\subsubsection{Pagine} Ogni pagina dev'essere numerata e deve riportare in testata il titolo del documento e della sezione che ospita. % e magari il nome/logo del gruppo e il nome del progetto

\subsubsection{Struttura} Dev'esserci un'unica struttura di sezioni, sottosezioni e paragrafi, comune a tutti i documenti dello stesso tipo.

\subsubsection{Indice dei contenuti} Documenti con più di tre o quattro sezioni devono riportare, all'inizio, un indice dei contenuti.

\subsubsection{Colori} Andrebbe minimizzato l'uso di colori, al fine di evitare ambiguità nel caso un documento venga stampato in bassa qualità.

\subsubsection{Glossario} All'interno di una sezione, la prima occorrenza di un termine che si trovi nel glossario dev'essere segnata con un pedice (tramite la macro \texttt{\textbackslash gloss\{termine\}}) in modo da far capire che si trova nel glossario.

\subsubsection{Strumenti} \dots

\subsubsection{Documenti finali} Andranno redatti i seguenti documenti.
\paragraph{Studio di fattibilità} Documento interno che vaglia i \emph{pro} e i \emph{contra} di ogni capitolato d'appalto.
\paragraph{Norme di progetto} Documento interno che norma i processi del progetto.
\paragraph{Piano di progetto} \dots
\paragraph{Piano di qualifica} \dots
\paragraph{Analisi dei requisiti} \dots
\paragraph{Specifica tecnica} \dots
\paragraph{Definizione di prodotto} \dots
\paragraph{Glossario} Documento esterno elencante specifici termini che il gruppo ha ritenuto opportuno definire. I motivi che portano ad inserire un termine nel glossario sono la sua potenziale ambiguità e/o la sua natura “tecnica”.
\paragraph{Manuale utente} \dots
\paragraph{Verbali} Documento esterno riportante i verbali delle riunioni del gruppo, sia quelle interne al gruppo sia quelle in presenza del committente.

\subsection{Processo di verifica} \dots

\subsection{Strumenti} \dots

\section{Processi organizzativi}

\subsection{Processo di miglioramento} % da ISO 12207-1995; pensavo fosse importante quindi l'ho aggiunto
Man mano che si procede nel progetto, è bene stare attenti a non lasciare che le attività ripetitive degradino in qualità. Ogni processo, essendo un insieme di attività che si ripetono, deve quindi abbracciare il modello \gloss{PDCA}. % come? 

\subsection{Strumenti}

\end{document}
